
\documentclass{article}
\usepackage{colortbl}
\usepackage{makecell}
\usepackage{multirow}
\usepackage{supertabular}

\begin{document}

\newcounter{utterance}

\centering \large Interaction Transcript for game `referencegame', experiment `number\_grids', episode 16 with reflexion{-}Qwen3.5{-}4B.
\vspace{24pt}

{ \footnotesize  \setcounter{utterance}{1}
\setlength{\tabcolsep}{0pt}
\begin{supertabular}{c@{$\;$}|p{.15\linewidth}@{}p{.15\linewidth}p{.15\linewidth}p{.15\linewidth}p{.15\linewidth}p{.15\linewidth}}
    \# & $\;$A & \multicolumn{4}{c}{Game Master} & $\;\:$B\\
    \hline

    \theutterance \stepcounter{utterance}  
    & & \multicolumn{4}{p{0.6\linewidth}}{
        \cellcolor[rgb]{0.9,0.9,0.9}{
            \makecell[{{p{\linewidth}}}]{
                \texttt{\tiny{[P1$\langle$GM]}}
                \texttt{You have attempted this same task before and did not succeed. These are the plans you wrote after those attempts.} \\
\texttt{Trial \#1: New plan: I failed because my description was too generic and failed to uniquely distinguish the target grid from the distractors. Specifically, I described the common features (solid bars in the first column and first/third rows) which are shared by the target and distractor 1, and I missed the critical difference: the target grid has an 'X' in the bottom{-}right corner (row 5, column 5), while distractor 1 has a '▢' there. Next time, I will construct the expression by first stating the unique feature that separates the target from the others (the bottom{-}right 'X'), then describing the shared structural elements (the solid bars) to confirm the overall shape, ensuring the description is logically sufficient to identify the target among the three options.} \\
\\ 
\\ 
\texttt{You are given three grids, where each of them is 5 by 5 in size.} \\
\texttt{Grids have empty cells marked with "▢" and filled cells marked with "X".} \\
\texttt{Your task is to generate a referring expression that best describes the target grid while distinguishing it from the two other distractor grids.} \\
\texttt{The first grid is the target grid, and the following two grids are the distractors.} \\
\\ 
\texttt{Target grid:} \\
\\ 
\texttt{X X X X □} \\
\texttt{X □ □ □ □} \\
\texttt{X X X X □} \\
\texttt{X □ □ □ X} \\
\texttt{X X X X □} \\
\\ 
\texttt{Distractor grid 1:} \\
\\ 
\texttt{X X X X □} \\
\texttt{X □ □ □ □} \\
\texttt{X X X X □} \\
\texttt{□ □ □ X □} \\
\texttt{X X X X □} \\
\\ 
\texttt{Distractor grid 2:} \\
\\ 
\texttt{□ X X X X} \\
\texttt{□ X □ □ X} \\
\texttt{□ X X X X} \\
\texttt{□ X □ □ X} \\
\texttt{□ X X X X} \\
\\ 
\texttt{Instruction: Describe the target grid.} \\
\texttt{Generate the referring expression starting with the tag "Expression: " for the given target grid. Omit any other text.} \\
            }
        }
    }
    & \\ \\

    \theutterance \stepcounter{utterance}  
    & \multicolumn{4}{p{0.6\linewidth}}{
        \cellcolor[rgb]{0.9,0.9,0.9}{
            \makecell[{{p{\linewidth}}}]{
                \texttt{\tiny{[P1$\rangle$GM]}}
                \texttt{Expression: The grid with a solid horizontal bar of five filled cells in the first row, a solid horizontal bar of five filled cells in the third row, a solid horizontal bar of five filled cells in the fifth row, and a vertical bar of five filled cells in the first column, with the bottom{-}right corner cell filled.} \\
            }
        }
    }
    & & \\ \\

    \theutterance \stepcounter{utterance}  
    & & & \multicolumn{2}{p{0.3\linewidth}}{
        \cellcolor[rgb]{0.9,0.9,0.9}{
            \makecell[{{p{\linewidth}}}]{
                \texttt{\tiny{[GM$|$GM]}}
                \texttt{Expression: The grid with a solid horizontal bar of five filled cells in the first row, a solid horizontal bar of five filled cells in the third row, a solid horizontal bar of five filled cells in the fifth row, and a vertical bar of five filled cells in the first column, with the bottom{-}right corner cell filled.} \\
            }
        }
    }
    & & \\ \\

    \theutterance \stepcounter{utterance}  
    & & \multicolumn{4}{p{0.6\linewidth}}{
        \cellcolor[rgb]{0.9,0.9,0.9}{
            \makecell[{{p{\linewidth}}}]{
                \texttt{\tiny{[GM$\rangle$P2]}}
                \texttt{You have attempted this same task before and did not succeed. These are the plans you wrote after those attempts.} \\
\texttt{Trial \#1: New plan: I incorrectly identified the target grid by relying on a partial match of the top row and first column, overlooking the specific requirement for a solid horizontal bar of four X's in the third row. The correct grid (Second) has X's in the third row at columns 1, 2, 3, and 4, whereas the First grid has an empty cell at column 4 in that row. Next time, I will systematically verify every constraint in the expression against the full grid pattern, specifically checking the third row for four consecutive X's and the bottom{-}right corner, rather than stopping after confirming the first two conditions.} \\
\\ 
\\ 
\texttt{You are given three grids, where each of them is 5 by 5 in size.} \\
\texttt{Grids have empty cells marked with "▢" and filled cells marked with "X".} \\
\texttt{You are also given a referring expression that describes one of the given grids.} \\
\texttt{Your task is to select the grid that matches the given referring expression.} \\
\texttt{Generate only the number (in text) of the grid that matches the given expression by selecting first, second, or third.} \\
\\ 
\texttt{First:} \\
\\ 
\texttt{X X X X □} \\
\texttt{X □ □ □ □} \\
\texttt{X X X X □} \\
\texttt{□ □ □ X □} \\
\texttt{X X X X □} \\
\\ 
\texttt{Second:} \\
\\ 
\texttt{X X X X □} \\
\texttt{X □ □ □ □} \\
\texttt{X X X X □} \\
\texttt{X □ □ □ X} \\
\texttt{X X X X □} \\
\\ 
\texttt{Third:} \\
\\ 
\texttt{□ X X X X} \\
\texttt{□ X □ □ X} \\
\texttt{□ X X X X} \\
\texttt{□ X □ □ X} \\
\texttt{□ X X X X} \\
\\ 
\texttt{Expression: The grid with a solid horizontal bar of five filled cells in the first row, a solid horizontal bar of five filled cells in the third row, a solid horizontal bar of five filled cells in the fifth row, and a vertical bar of five filled cells in the first column, with the bottom{-}right corner cell filled.} \\
\texttt{Question: Which grid does the expression refer to?} \\
\texttt{Start with the tag "Answer: ", followed by your selection. Omit any other text.} \\
            }
        }
    }
    & \\ \\

    \theutterance \stepcounter{utterance}  
    & & & \multicolumn{4}{p{0.6\linewidth}}{
        \cellcolor[rgb]{0.9,0.9,0.9}{
            \makecell[{{p{\linewidth}}}]{
                \texttt{\tiny{[GM$\langle$P2]}}
                \texttt{Answer: Third} \\
            }
        }
    }
     \\ \\

    \theutterance \stepcounter{utterance}  
    & & & \multicolumn{2}{p{0.3\linewidth}}{
        \cellcolor[rgb]{0.9,0.9,0.9}{
            \makecell[{{p{\linewidth}}}]{
                \texttt{\tiny{[GM$|$GM]}}
                \texttt{third} \\
            }
        }
    }
    & & \\ \\

    \theutterance \stepcounter{utterance}  
    & \multicolumn{4}{p{0.6\linewidth}}{
        \cellcolor[rgb]{0.95,0.95,0.95}{
            \makecell[{{p{\linewidth}}}]{
                \texttt{\tiny{[P2$\rightarrow$P1]}}
                \texttt{Post{-}episode Reflexion (not sent to GM):} \\
\texttt{New plan: I failed because my description listed the shared structural elements (solid bars) first, which made the target grid indistinguishable from Distractor 1. The critical differentiator is the bottom{-}right cell: the target has an 'X' while Distractor 1 has a '▢'. Next time, I will start the expression by explicitly stating the unique bottom{-}right 'X' to immediately separate the target from Distractor 1, and then describe the common row/column patterns to confirm the overall shape against Distractor 2.} \\
            }
        }
    }
    & & \\ \\

    \theutterance \stepcounter{utterance}  
    & & & \multicolumn{4}{p{0.6\linewidth}}{
        \cellcolor[rgb]{0.95,0.95,0.95}{
            \makecell[{{p{\linewidth}}}]{
                \texttt{\tiny{[P1$\rightarrow$P2]}}
                \texttt{Post{-}episode Reflexion (not sent to GM):} \\
\texttt{New plan: I incorrectly identified the target grid by focusing on the first row and first column, which were identical across all three grids, and failing to verify the specific constraints for the third and fifth rows. The expression requires solid horizontal bars in the third and fifth rows, but the First grid has empty cells in those rows, and the Second grid has a filled cell in the bottom{-}right corner instead of an empty one. Next time, I will systematically check every row and column constraint in the expression against the full grid pattern, specifically verifying that the third and fifth rows are completely filled with X's and that the bottom{-}right corner is empty, rather than relying on partial matches.} \\
            }
        }
    }
     \\ \\

\end{supertabular}
}

\end{document}
